% Shared titles: update each definition to change both the TOC and first slide title.
\newcommand{\lossestitle}{The Night Asgard Lost Its King}
\newcommand{\sakaartitle}{Cast Out by the Bifrost}
\newcommand{\decisionstitle}{The Price of the Road Home}
\newcommand{\prioritiestitle}{Asgard Is Its People}
\newcommand{\returntitle}{The Road Home Ends in Fire}
\newcommand{\modeltitle}{A Tally of the Survivors}

{
    \setbeamertemplate{footline}{}
    \begin{frame}[plain,noframenumbering]
        \titlepage
        \note{Introduce the story as a journey from Asgard's fall to the survival of its people.}
    \end{frame}
}

\begin{frame}{\contentslabel}{From the fall of Asgard to the flight of its people}
    \begin{SlideContents}
    \end{SlideContents}
    \note{Preview the move from exile to the Asgardians' escape.}
\end{frame}

\section{The Fall and Exile}
\subsection{\lossestitle}
\begin{frame}{\lossestitle}{Odin's death, Mjölnir's fall, and Hela's return}
    % The image is scaled to its specified height and cropped to the column width.
    \begin{SlideBulletsPhoto}{assets/hela_mjolnir.jpg}{4.2cm}{Hela catches Mjölnir.}
        \begin{itemize}
            \item \textbf{Odin} is gone; the Allfather's reign has ended.
            \item \textbf{Mjölnir} lies shattered, and Thor's certainty with it.
            \item \textbf{Hela} claims the throne and casts Asgard into shadow.
        \end{itemize}
    \end{SlideBulletsPhoto}
    \note{Thor loses Odin, Mjölnir, and his place in Asgard.}
\end{frame}

\subsection{\sakaartitle}
\begin{frame}{\sakaartitle}{Thor's fall into the Grandmaster's arena}
    \begin{SlideTitlePhoto}{assets/thor_hulk_arena.jpg}{5.2cm}{0.95\linewidth}
        % Scale to the specified height and crop any horizontal overflow.
        \caption{Thor faces the Hulk in Sakaar's arena.}
    \end{SlideTitlePhoto}
    \note{Thor's exile on Sakaar gives him a new path back.}
\end{frame}

\begin{frame}{A Kingdom Can Fall. Its People Can Live.}{A kingdom and its people are not the same thing}
    \begin{SlidePhotoPair}
        % Both paired images use the same height and width limits.
        \SlidePhotoPanel{assets/thor_sakaar.jpg}{3.8cm}{Thor surveys Sakaar.}
        \hfill
        \SlidePhotoPanel{assets/thor_loki_alliance.jpg}{3.8cm}{Loki's uneasy alliance opens a path home.}
        \caption{On Sakaar, the brothers find a way back to the people under Hela's rule.}
    \end{SlidePhotoPair}
    \note{The brothers turn exile into a chance to help their people.}
\end{frame}

\subsubsection{Brothers at a Crossroads}
\begin{frame}{Brothers at a Crossroads}{Old betrayals shadow their alliance}
    \begin{SlideTitlePhoto}{assets/thor_loki_elevator.jpg}{5.2cm}{0.95\linewidth}
        % Scale to the specified height and crop any horizontal overflow.
        \caption{Thor knows Loki may still choose his own escape.}
    \end{SlideTitlePhoto}
    \note{Loki begins to choose alliance over his own escape.}
\end{frame}

\subsection{\decisionstitle}
\begin{frame}{\decisionstitle}{Three losses and the choice to return}
    \begin{SlideThreeColumns}
        \begin{column}{0.30\textwidth}
            \textbf{What Was Lost}
            \begin{itemize}
                \item Odin, his father
                \item Mjölnir, his weapon
                \item Asgard under Hela's rule
            \end{itemize}
        \end{column}
        \begin{column}{0.30\textwidth}
            \textbf{The Road Back}
            \begin{enumerate}
                \item Survive Sakaar
                \item Find Loki again
                \item Return to Asgard
            \end{enumerate}
        \end{column}
        \begin{column}{0.30\textwidth}
            \textbf{The Choice}
            \begin{description}
                \item[Thor] \keyphrase{Save the people}
                \item[Loki] Return with them
                \item[People] Escape Asgard
            \end{description}
        \end{column}
    \end{SlideThreeColumns}
    \note{Thor chooses to save Asgard's people even if the kingdom is lost.}
\end{frame}

\begin{frame}{The Road to Exodus}{From Asgard's fall to the survivors' escape}
    \SlideTimeline
    \begin{tikzpicture}[
            stage/.style={draw=softblack,align=center,minimum width=3cm,minimum height=1cm},
            every node/.style={font=\small}]
        \node[stage] (loss) {Fall\\Asgard};
        \node[stage,right of=loss,node distance=4cm] (exile) {Exile\\Sakaar};
        \node[stage,final stage,right of=exile,node distance=4cm] (return) {Ragnarok\\The people escape};
        \draw[progress arrow] (loss) -- (exile);
        \draw[progress arrow] (exile) -- (return);
    \end{tikzpicture}
    \note{Trace the journey from Asgard through Sakaar to the escape.}
\end{frame}

\section{Thor's Choice}
{
    \setbeamertemplate{frametitle}{}
    \begin{frame}
        \begin{tikzpicture}[remember picture,overlay]
            \node[inner sep=0pt,font=\LARGE] at (current page.center) {The people come first.};
        \end{tikzpicture}
        \note{Thor chooses to put the Asgardians first.}
    \end{frame}
}

\subsection{\prioritiestitle}
\begin{frame}{\prioritiestitle}
    \centering
    {\Large Asgard is not its golden halls. \keyphrase{It is its people.}\par}
    \medskip
    {\footnotesize From the royal archives: \citet{lore_thor,lore_loki,lore_asgard}.\par}
    \note{Emphasize that Asgard lives through its people.}
\end{frame}

\begin{frame}{A King Without a Kingdom}{Thor learns that Asgard lives on in its people}
    \begin{SlideTable}
        \caption{The path from Hela's return to the people's escape.}
        \setlength\extrarowheight{4pt}
        \begin{tabularx}{\linewidth}{lYY}
            \toprule
            \rowcolor{offwhite}
            \textbf{Stage} & \textbf{Disruption} & \textbf{Response}\\
            \midrule
            Asgard & Odin is dead; Mjölnir is shattered; Hela returns & Face the fall\\
            Sakaar & Exile in the Grandmaster's arena & Gather allies\\
            Return & Hela rules Asgard & Bring the people to safety\\
            \midrule
            \highlightrow
            Resolution & Asgard falls to Ragnarok & Carry its people into the stars\\
            \bottomrule
        \end{tabularx}
    \end{SlideTable}
    \note{Follow how each setback leads toward the people's escape.}
\end{frame}

\subsection{\returntitle}
\begin{frame}{\returntitle}{The choice that saves the Asgardians}
    \begin{block}{The Loss}
        Odin is gone. Mjölnir is shattered. Hela rules Asgard.
    \end{block}
    \begin{alertblock}{The Alliance}
        On Sakaar, Thor joins forces with Loki, Valkyrie, and the Hulk.
    \end{alertblock}
    \begin{exampleblock}{The Choice}
        Thor sends Loki to awaken Surtur, accepting Asgard's fall so its people can escape.\footnote{A chronicle of Asgard's fall and the people's escape.}
    \end{exampleblock}
    \medskip
    From the royal archives: \citet{lore_thor}.
    \note{Thor's choice gives the Asgardians a chance to escape.}
\end{frame}

\subsection{\modeltitle}
\begin{frame}[fragile]{\modeltitle}{An imagined count of those who escaped}
    \begin{SlideCodeExplanation}
        \begin{column}{0.46\textwidth}
            \begin{Verbatim}[fontsize=\footnotesize]
def protected_fraction(p, N):
    if N <= 0 or p < 0 or p > N:
        raise ValueError("Bad input")
    return p / N

saved, total = 75, 100
fraction = protected_fraction(saved, total)
print(f"Saved: {fraction:.0%}")
      \end{Verbatim}
        \end{column}
        \begin{column}{0.46\textwidth}
            \textbf{This imagined tally} records the Asgardians who reach safety.

            \medskip
            This is an imagined count, not a number given in the chronicle:
            \[
                r=\frac{p}{N},\qquad N>0,\quad 0\leq p\leq N.
            \]
        \end{column}
    \end{SlideCodeExplanation}
    \note{The 75 percent figure is an imagined example, not a historical count.}
\end{frame}

\begin{frame}{A Narrowing Window to Escape}{A hypothetical account of the flight from Asgard}
    This imagined model estimates the chance of getting the Asgardians to safety:
    \begin{equation}
        P(E)=\frac{1}{1+\exp[-(\alpha C-\gamma D)]},\qquad \alpha,\gamma>0.
    \end{equation}
    Here, $C$ is available transport capacity and $D$ is the time lost before the escape begins.
    This model is imagined; the chronicle gives no exact evacuation data.

    \medskip
    Greater transport capacity and an earlier escape give more Asgardians a chance to survive Surtur's fire.
    \note{More transport and less delay improve the chance of escape in this imagined model.}
\end{frame}

\section{\referenceslabel}
\setbeamertemplate{frametitle continuation}[from second]
\begin{frame}[allowframebreaks]{\referenceslabel}{Records of Asgard's fall and the survivors' flight}
    \begin{SlideReferences}
    \bibliographystyle{plainnat}
    \bibliography{references}
    \end{SlideReferences}
    \note{Point to the sources listed here for the chronicle.}
\end{frame}
